% Tabelas e frases de apoio aos Resultados do artigo (docs/article/artigo.tex).
% Saíram do artigo, palavra por palavra, na simplificação das métricas; os
% números continuam conferidos contra os CSV por scripts/check-article-numbers.py.
% Documento independente: NÃO é incluído no artigo por \input.
\documentclass[12pt]{article}
\usepackage[utf8]{inputenc}
\usepackage[T1]{fontenc}
\usepackage{booktabs}
\usepackage[brazil]{babel}

\begin{document}

\section*{Tabelas e detalhes de apoio aos Resultados}

As tabelas e frases abaixo estavam no artigo e foram retiradas para encurtá-lo. Os números são os mesmos do artigo e são conferidos contra os resultados brutos por \texttt{scripts/check-article-numbers.py}.

\subsection*{Busca em C1 por tipo de consulta (Seção 8.1)}

\textit{Entre os 5 textos.} As duas buscas por palavras acertam as 5 consultas literais; nas 7 paráfrases, a contagem acerta 4 (57,1\%) e o BM25, 5 (71,4\%). A busca densa acerta 4 das 7 paráfrases, mas só 3 das 5 literais (60,0\%): quase não depende das palavras usadas, mas erra consultas que repetem o texto (Anexo~A7, Tabela~2).

\textit{Entre as 19 seções.} A dependência das palavras também cresce: a contagem continua acertando as 5 literais, mas só 1 das 7 paráfrases (14,3\%), e o BM25 acerta 3 das 7.

\subsection*{Origem das fontes no índice (Seção 8.4)}

\begin{table}[!htbp]
\centering
\caption{Unidades de C1 (155) com origem completa no banco, a cada correção.}
\label{apoio:provenance-evolucao}
\begin{tabular}{lccc}
\toprule
 & Critério original & + página no trecho & + origem declarada \\
\midrule
Unidades completas & 51 & 70 & 155 \\
Origem completa (\%) & 32,9 & 45,2 & 100,0 \\
\bottomrule
\end{tabular}
\end{table}

\subsection*{Trechos recuperados contra documento inteiro (Seção 8.5)}

\begin{table}[!htbp]
\centering
\caption{Acervo de ensino (desenvolvimento), 72 perguntas, três repetições: porcentagem de respostas certas e de respostas que citam a página certa no texto, com o sistema buscando trechos e com o PDF inteiro entregue ao modelo. A página citada conta as 201 respostas que têm página-alvo.}
\label{apoio:geracao}
\begin{tabular}{lrrrr}
\toprule
 & \multicolumn{2}{c}{Acerto (\%)} & \multicolumn{2}{c}{Página citada certa (\%)} \\
\cmidrule(lr){2-3}\cmidrule(lr){4-5}
Tipo de pergunta & Trechos & PDF inteiro & Trechos & PDF inteiro \\
\midrule
Prosa (30)      & 100,0 & 100,0 & 93,3 & 88,9 \\
Tabela (21)     & 100,0 & 100,0 & 90,5 & 96,8 \\
Calendário (16) & 100,0 & 100,0 & 100,0 & 100,0 \\
Negativa (5)    & 100,0 & 80,0 & -- & -- \\
\midrule
Todas (72)      & 100,0 & 98,6 & 94,0 & 94,0 \\
\bottomrule
\end{tabular}
\end{table}

\textit{Custo.} A razão entre os \emph{tokens} de F e de T, cuja mediana foi 26,5, variou de 0{,}7 a 96 vezes. Em média, T enviou 4.215 \emph{tokens} por pergunta e F, 72.067; sem as negativas, em que F recebe cinco documentos, foram 4.288 e 58.751. Na execução exploratória com outro modelo, de janela de 131.072 \emph{tokens}, o provedor recusou o \emph{prompt} de F nas cinco negativas, com cinco documentos (de 227.647 a 227.653 \emph{tokens}).

\subsection*{Legislação com norma revogada (Seção 8.6)}

\begin{table}[!htbp]
\centering
\caption{Acervo de legislação: 19 perguntas, uma repetição. A página certa entre as fontes conta as 16 perguntas que têm página-alvo.}
\label{apoio:dominio2}
\begin{tabular}{lrrr}
\toprule
Tipo de pergunta & Perguntas & Acerto (\%) & Página certa entre as fontes (\%) \\
\midrule
Divergente & 5 & 60,0 & 80,0 \\
Temporal   & 3 & 100,0 & 100,0 \\
Prosa      & 8 & 75,0 & 75,0 \\
Negativa   & 3 & 100,0 & --    \\
\midrule
Total      & 19 & 78,9 & 81,2 \\
\bottomrule
\end{tabular}
\end{table}

\subsection*{Bateria em três domínios (Seção 8.7)}

\begin{table}[!htbp]
\centering
\caption{Bateria, pergunta original: acerto por acervo, em porcentagem das respostas (cada pergunta foi feita três vezes), com o intervalo de confiança de 95\%, calculado por \emph{bootstrap}; a última coluna conta as perguntas que acertaram nas três vezes.}
\label{apoio:bateria-tema}
\begin{tabular}{lrrr}
\toprule
Acervo & Acerto (\%) & Intervalo de 95\% & Certas nas 3 vezes \\
\midrule
Segurança do trabalho & 62,3 & 43,5--81,2 & 14 de 23 \\
Produto e consumidor  & 82,6 & 65,2--95,7 & 19 de 23 \\
Licitações            & 64,0 & 44,0--84,0 & 16 de 25 \\
\midrule
Total                 & 69,5 & 59,2--79,3 & 49 de 71 \\
\bottomrule
\end{tabular}
\end{table}

\begin{table}[!htbp]
\centering
\caption{Bateria, pergunta original: acerto por forma da pergunta, nos três acervos (respostas certas sobre o total).}
\label{apoio:bateria-forma}
\begin{tabular}{lrr}
\toprule
Forma & Acertos & Acerto (\%) \\
\midrule
Negativa       & 27/27 & 100,0 \\
Existência     & 33/36 & 91,7 \\
Confirmação    & 27/36 & 75,0 \\
Aberta         & 48/78 & 61,5 \\
Premissa falsa & 13/36 & 36,1 \\
\bottomrule
\end{tabular}
\end{table}

Na forma existência, cuja resposta certa é sempre ``sim'', responder ``Sim.'' a tudo acertaria as 12 perguntas, de modo que o 91,7\% dessa forma só mostra que o sistema não negou um assunto presente.

\subsection*{Expansão da pergunta (Seção 8.9)}

Linhas retiradas da tabela da expansão (sem expansão, com expansão):

\begin{tabular}{lrr}
\toprule
 & Sem expansão & Com expansão \\
\midrule
Página certa entre as fontes (\%)                  & 58,1 & 82,3 \\
\emph{Tokens} enviados por pergunta                & 4.699 & 5.327 \\
\bottomrule
\end{tabular}

Contando por pergunta, o ganho foi de 19,7 pontos percentuais (de 49 para 63 das 71 perguntas certas nas três repetições), com intervalo de 95\% de 10,5 a 29,0 pontos, inteiro acima de zero. Por resposta, os dois intervalos nem se encostam: com expansão, o acerto fica no intervalo de 80,3\% a 95,8\%, contra 59,2\% a 79,3\% sem ela. Das perguntas que erravam sem expansão, a expansão corrigiu 13 de 18 paráfrases e 1 de 4 literais, na direção prevista.

\end{document}
